% chapters/testing/strategy.tex
% Sección 6.1: Estrategia de pruebas (~1 página)
% ============================================================================
%
% GUÍA: Describir la estrategia global de pruebas.
%
%   1. Niveles de prueba aplicados (pirámide de testing):
%      - Pruebas unitarias: lógica de negocio aislada.
%      - Pruebas de integración: endpoints completos con BD.
%      - Pruebas de carga: rendimiento bajo estrés.
%      - Pruebas de aceptación: verificación de requisitos.
%
%   2. Herramientas:
%      - pytest + pytest-django (unitarias e integración).
%      - factory_boy o model_bakery (fixtures/factories).
%      - TODO: Especificar herramienta de pruebas de carga
%        (Locust, k6, Artillery, ApacheBench, etc.).
%      - Cobertura: coverage.py / pytest-cov.
%
%   3. Entorno de pruebas:
%      - Base de datos PostgreSQL de test (transaccional).
%      - MinIO de test (bucket separado).
%      - Redis de test o mock.
%
%   4. Criterios de aceptación:
%      - Cobertura mínima objetivo (p.ej., 80% de la lógica de negocio).
%      - Todos los endpoints principales cubiertos por integración.
%      - Tiempos de respuesta aceptables bajo carga.
%
%   Opcionalmente, incluir una figura con la pirámide de testing:
%   \begin{figure}{fig:test_pyramid}{Pirámide de pruebas del proyecto}
%     \image{}{}{test_pyramid}
%   \end{figure}
%   TODO: Crear figures/test_pyramid.png (opcional, sencillo).
% ============================================================================

% TODO: Redactar la estrategia de pruebas.
